% ============================================================
%  Liste des Abréviations — Mémoire de Master BIBDA
%  CHOUFA Zouhair
%  NOTE : nécessite \usepackage{longtable} dans le préambule
%         de main.tex (voir instructions fournies).
% ============================================================
\chapter*{Liste des Abréviations}
\addcontentsline{toc}{chapter}{Liste des Abréviations}

\begin{longtable}{p{3.2cm} p{10.8cm}}
\hline
\textbf{Abréviation} & \textbf{Signification} \\
\hline
\endfirsthead
\multicolumn{2}{l}{\textit{(suite de la page précédente)}} \\
\hline
\textbf{Abréviation} & \textbf{Signification} \\
\hline
\endhead
\hline
\multicolumn{2}{r}{\textit{(suite sur la page suivante)}} \\
\endfoot
\hline
\endlastfoot

API       & Application Programming Interface (interface de programmation applicative) \\
BD        & Big Data \\
BIBDA     & Business Intelligence and Big Data Analytics (intitulé du Master) \\
CLI       & Command Line Interface (interface en ligne de commande) \\
CRM       & Customer Relationship Management (gestion de la relation client) \\
CSV       & Comma-Separated Values (valeurs séparées par des virgules) \\
DL        & Deep Learning (apprentissage profond) \\
EM        & Expectation-Maximization (algorithme espérance-maximisation) \\
FOAF      & Friend Of A Friend (vocabulaire RDF de description de réseaux sociaux) \\
GPT       & Generative Pre-trained Transformer \\
GraphDB   & Triplestore RDF développé par Ontotext, utilisé pour le stockage du KG produit \\
IA        & Intelligence Artificielle \\
ISO       & International Organization for Standardization (organisation internationale de normalisation) \\
JSON      & JavaScript Object Notation (format d'échange de données) \\
KG        & Knowledge Graph (graphe de connaissances) \\
KGs       & Knowledge Graphs (pluriel) \\
LangChain & Framework open source d'orchestration d'agents et de \textit{prompts} pour LLMs \\
LLaMA     & Large Language Model Meta AI (famille de modèles de langage de Meta) \\
LLM       & Large Language Model (grand modèle de langage) \\
LLMs      & Large Language Models (pluriel) \\
LLMatch   & Méthode de \textit{schema matching} fondée sur les LLMs (littérature) \\
Morph-KGC & Moteur de triplification déclarative RML/YARRRML utilisé au Module~3 du pipeline \\
NLP       & Natural Language Processing (traitement automatique du langage naturel) \\
OAEI      & Ontology Alignment Evaluation Initiative (initiative d'évaluation de l'alignement d'ontologies) \\
OWL       & Web Ontology Language (langage d'ontologies du Web sémantique) \\
PySHACL   & Bibliothèque Python d'implémentation du langage de contraintes SHACL \\
PyVis     & Bibliothèque Python de visualisation interactive de graphes \\
R2RML     & RDB to RDF Mapping Language (standard W3C de mapping base relationnelle vers RDF) \\
RDB       & Relational DataBase (base de données relationnelle) \\
RDF       & Resource Description Framework (modèle de données du Web sémantique) \\
RDFLib    & Bibliothèque Python de manipulation de graphes RDF \\
REST      & Representational State Transfer (style d'architecture d'API Web) \\
RML       & RDF Mapping Language (standard W3C de mapping générique vers RDF) \\
SBERT     & Sentence-BERT (modèle d'embeddings de phrases fondé sur BERT) \\
SCHEMORA  & Méthode d'alignement de schémas par \textit{prompt engineering} (littérature) \\
SemTab    & Semantic Table Interpretation Challenge (banc d'essai communautaire d'interprétation sémantique de tables) \\
SHACL     & Shapes Constraint Language (langage W3C de validation de graphes RDF) \\
SIRH      & Système d'Information des Ressources Humaines \\
SPARQL    & SPARQL Protocol And RDF Query Language (langage de requête pour RDF) \\
SQL       & Structured Query Language (langage de requête structuré) \\
SQLite    & Système de gestion de base de données relationnelle embarqué \\
STI       & Semantic Table Interpretation (interprétation sémantique de tables) \\
TF-IDF    & Term Frequency--Inverse Document Frequency (pondération statistique de termes) \\
URI       & Uniform Resource Identifier (identifiant uniforme de ressource) \\
URIs      & Uniform Resource Identifiers (pluriel) \\
UTF-8     & Unicode Transformation Format -- 8 bits (encodage de caractères) \\
W3C       & World Wide Web Consortium (organisme de standardisation du Web) \\
XML       & eXtensible Markup Language (langage de balisage extensible) \\
XSD       & XML Schema Definition (langage de définition de schémas XML) \\
YAML      & YAML Ain't Markup Language (format de sérialisation de données lisible) \\
YARRRML   & Yet Another RML Rule ML (langage déclaratif simplifié pour la génération de règles RML) \\
$\tau$    & Seuil de décision du Module~2 (schema matching hybride) \\
$\lambda$ & Seuil de décision du Module~4 (résolution d'entités, modèle de Fellegi-Sunter) \\

\end{longtable}

\newpage